\documentclass[a4paper,11pt]{article}

\usepackage[margin=2.5cm]{geometry}
\usepackage[T1]{fontenc}
\usepackage[utf8]{inputenc}
\usepackage{lmodern}
\usepackage{booktabs}
\usepackage{tabularx}
\usepackage{array}
\usepackage{enumitem}
\usepackage{hyperref}

\hypersetup{
  colorlinks=true,
  linkcolor=black,
  urlcolor=blue,
  citecolor=black
}

\setlength{\parindent}{0pt}
\setlength{\parskip}{0.45em}
\renewcommand{\arraystretch}{1.15}
\newcolumntype{Y}{>{\raggedright\arraybackslash}X}

\title{\textbf{Benchmarking LLM-Based Honeypots with a Scripted Local Harness: Setup, Findings and Limitations}}
\author{Muhammad Rumman Aslam}
\date{}

\begin{document}
\maketitle

\begin{center}
\textbf{Final batch, collected Nov 2026, 5 runs per target, 120\,s timeout.
Supersedes the pilot batch.}
\end{center}

\begin{abstract}
This report documents a scripted local benchmark of LLM-based SSH and HTTP
honeypots, a static Cowrie baseline, ground-truth services, an ANSI-hidden
prompt-injection detector target, and HoneyBee as a separate environment
generator. The final batch used five runs per SSH and HTTP target. Each SSH
run sent 37 ordered commands, of which 36 were scored; each HTTP run sent 20
requests. The harness used a 120-second runner timeout and recorded raw JSONL
observations for later analysis. SSH latency summaries use completed commands
only, while timed-out observations are retained as censored events. Similarity
comparisons use ANSI-stripped prefixes capped at 4,000 characters. The
observed data show target-specific determinism, command-specific timeouts,
large outputs, and differing HTTP response consistency. HoneyBee was evaluated
as an LLM generator rather than as a honeypot. These are pilot-scale,
scripted measurements from one shared model and endpoint; they do not rank
tools, establish causes, or generalize beyond the tested configurations.
\end{abstract}

\section{Introduction and goals}

LLM-based honeypots can answer shell or HTTP interactions with outputs
constructed by a language model rather than by a fixed emulation layer. A
benchmark must therefore record both conventional service observations and
properties specific to generated text: repeat consistency, output length,
interior prompt leakage, empty or unexpected outputs, and timeout behavior.
This project uses a local harness to make those observations comparable across
several SSH and HTTP targets.

The study is pilot-scale. It uses a single shared model and API endpoint for
the LLM targets, five scripted runs, and a fixed command/request corpus. The
goal is measurement and reproducibility rather than a universal ranking. The
ground-truth SSH service is Ubuntu 20.04 with OpenSSH, and the HTTP reference
is Ubuntu 20.04 with nginx. Cowrie supplies a static SSH baseline. The
llm-honeypot target is treated as a detector target and excluded from fidelity,
similarity, determinism, and aggregate averages.

The report has three practical goals: document the local setup and target
configuration; state the metrics and accounting rules precisely; and preserve
the limitations and unresolved questions that should guide a larger study.

\section{Setup}

\subsection{Targets and ports}

All benchmark services were published on loopback ports. Table~\ref{tab:targets}
summarizes the final target inventory. HoneyBee is deliberately listed as a
generator: it produces Docker-oriented artifacts and is not an interactive
honeypot in this benchmark.

\begin{table}[ht]
\centering
\caption{Final target inventory.}
\label{tab:targets}
\small
\begin{tabularx}{\textwidth}{l l r l Y}
\toprule
Target & Track & Port & Role & Recorded scope \\
\midrule
DECEIVE & SSH & 8022 & LLM honeypot & Scripted SSH responses \\
Beelzebub & SSH & 2222 & LLM honeypot & Scripted SSH responses \\
Cowrie & SSH & 2223 & Static baseline & Static emulation \\
AdvancedshelLM & SSH & 2224 & LLM honeypot & Subprocess/wrapper path \\
Ground truth SSH & SSH & 2225 & Reference & Ubuntu 20.04/OpenSSH \\
llm-honeypot & SSH & 2227 & Detector target & ANSI-hidden marker detector \\
VelLMes & SSH & 2228 & LLM honeypot & Scripted SSH responses \\
galah & HTTP & 8088 & LLM HTTP responder & Scripted HTTP responses \\
Ground truth HTTP & HTTP & 2226 & Reference & Ubuntu 20.04/nginx \\
HoneyBee & generator & 2229 config & LLM generator & Output-quality evaluation \\
\bottomrule
\end{tabularx}
\end{table}

\subsection{Runtime and configuration}

The documented reproduction environment is Windows 10/11 with Python 3.10+
(tested with Python 3.12), Git, and Docker Desktop with WSL 2. Services are
bound to loopback. The LLM targets use one shared model/API endpoint, while
the configuration values that were explicitly recorded differ by target.
Table~\ref{tab:config} preserves ``not recorded'' where the permitted
documents did not specify a value.

\begin{table}[ht]
\centering
\caption{Recorded target configuration.}
\label{tab:config}
\scriptsize
\begin{tabularx}{\textwidth}{l Y l l l Y}
\toprule
Target & Stack & Model & Temperature & Max tokens & User/hostname and main configuration \\
\midrule
DECEIVE & Python 3.11-slim + DECEIVE & glm-5p3-flash & 0.2 & no output cap & guest / not recorded; output stripping, Fireworks environment, exit sentinel, lazy imports \\
Beelzebub & Go image + mounted YAML & glm-5p3-flash & not exposed by plugin config & not exposed by plugin config & guest / ubuntu; API-key/host fallback and benchmark YAML \\
AdvancedshelLM & Python 3.11-slim + shelLM & glm-5p3-flash & 0.2 & 4096 & guest / not recorded; client, exit handling, environment fallback, wrapper \\
VelLMes & Python 3.11-slim + VelLMes & glm-5p3-flash & 0.2 & 800 & user / not recorded; SSH wrapper and benchmark config \\
galah & Go 1.22 builder + Debian slim & glm-5p3-flash & 0.2 & not recorded & not applicable / not recorded; response cleaning and entrypoint/config \\
Cowrie & Official image digest & not applicable & not applicable & not applicable & root / ubuntu; mounted config, no upstream source edit \\
llm-honeypot & Cowrie v2.6.1 overlay & not applicable & not applicable & not applicable & root / not recorded; overlay Dockerfile and v2.6.1 pin \\
Ground truth SSH & Ubuntu 20.04 + OpenSSH & not applicable & not applicable & not applicable & guest / ubuntu; PTY and PAM benchmark changes \\
Ground truth HTTP & Ubuntu 20.04 + nginx 1.18.0 & not applicable & not applicable & not applicable & not applicable / not recorded; benchmark Dockerfile/compose \\
HoneyBee & Python 3.11-slim + Streamlit & glm-5p3-flash & 0.2 & 8192 & not applicable; deployment disabled and empty-generation handling \\
\bottomrule
\end{tabularx}
\end{table}

Third-party source changes are documented as patches and notes. The verified
third-party commit hashes are the HoneyBee clone
\texttt{dbf62c42604eeae2cfbfd4d36b99bef65555c303} and the VelLMes clone
\texttt{192fa5084df84b38a16f6581939fd3143c350181}. Other clone commit hashes
are omitted here. The ground-truth services are benchmark-created services,
not third-party source repositories.

HoneyBee is not a honeypot in this study. It is a Streamlit application that
uses an LLM to generate misconfigured Docker artifacts. The one-click deploy
path is disabled, the image has no Docker CLI or Docker socket mount, and the
evaluation reports generator validity rather than a honeypot security
assessment.

\section{Experimental protocol}

\subsection{Inputs and collection order}

The SSH manifest contains 37 ordered commands: 36 scored commands plus
\texttt{exit}. Its observed categories are recon, files, process, errors,
permissions, and state. Host-specific commands are excluded from fidelity
scoring. The HTTP manifest contains 20 requests: 17 GET, one HEAD, one
OPTIONS, and one POST, including three repeated \texttt{GET /} requests.

Five runs were collected for every SSH and HTTP target. The runner timeout was
120 seconds, and a duplicate guard skipped complete target/run pairs. The
raw-file timestamps show a two-pass SSH order. In run 1, ground-truth SSH,
Cowrie, llm-honeypot, and Beelzebub were written first, followed by
AdvancedshelLM, DECEIVE, and VelLMes. The same two-stage SSH order repeated
for runs 2--5. Galah HTTP runs 1--5 were collected after the first SSH pass;
ground-truth HTTP was written immediately after Galah run 5. Containers were
freshly started before the batch.

\subsection{Metric definitions}

\begin{itemize}[leftmargin=*]
\item \textbf{Latency}: command or request wall time for completed commands
only. SSH timeouts are censored at the runner timeout and excluded from
median and p95, so the resulting latency summaries understate latency for
targets with timeouts.
\item \textbf{Fidelity}: the mean \texttt{SequenceMatcher} ratio for a target
command against the first loaded ground-truth SSH output for that command.
Host-specific commands are reference-only. Before each comparison, ANSI
escapes are removed and both texts are capped at their first 4,000 characters.
\item \textbf{Determinism}: mean pairwise \texttt{SequenceMatcher} similarity
across non-empty run outputs for the same target and command, with the same
ANSI-stripped 4,000-character cap.
\item \textbf{Shared commands}: the 26 non-host-specific commands with a
non-timeout observation for every compared target. Shared fidelity averages
command means over this intersection.
\item \textbf{Oversized output}: a command output longer than 20,000
characters; the reported value counts such outputs per target.
\item \textbf{Transcript leakage}: an output containing an interior shell
prompt matching
\texttt{(?m)\textasciicircum\textbackslash s*[\textbackslash w.-]+@[\textbackslash w.-]+:[\textbackslash n]*[\$\#]\textbackslash s}.
The legitimate final non-empty prompt is excluded.
\end{itemize}

\section{Results}

\subsection{SSH results}

The scored SSH denominator is 180 commands per target. Table~\ref{tab:ssh}
reports completed observations separately from timeouts. The fidelity interval
is the shared-command interval over 26 commands. The detector target is shown
for accounting and leakage, but its fidelity and determinism are excluded.

\begin{table}[ht]
\centering
\caption{SSH results. Latency is median/p95 seconds over completed commands.}
\label{tab:ssh}
\scriptsize
\begin{tabularx}{\textwidth}{l r r r r l r l l}
\toprule
Target & Median/p95 & Completed & Timeouts & Empty & Fidelity (95\% CI) & Det. & Oversized & Leakage \\
\midrule
Ground truth & 0.562/0.579 & 180 & 0 & 35 & --- & --- & 0 & 0/180 \\
Cowrie & 0.563/0.610 & 180 & 0 & 40 & 0.5877 (0.4203--0.7438) & 0.9874 & 0 & 0/180 \\
AdvancedshelLM & 1.922/4.504 & 180 & 0 & 26 & 0.5947 (0.4336--0.7386) & 0.7836 & 0 & 0/180 \\
VelLMes & 1.687/4.874 & 171 & 9 & 25 & 0.6070 (0.4601--0.7506) & 0.8364 & 0 & 0/180 \\
DECEIVE & 1.984/10.361 & 175 & 5 & 27 & 0.4731 (0.3493--0.5858) & 0.4664 & 2 & 0/180 \\
Beelzebub & 5.023/31.656 & 164 & 16 & 16 & 0.3737 (0.2319--0.5239) & 0.6800 & 9 & 14/180 \\
llm-honeypot & detector & 180 & 0 & 30 & excluded & excluded & 0 & 0/180 \\
\bottomrule
\end{tabularx}
\end{table}

Timeouts were command-specific. VelLMes \texttt{ps\_aux} timed out in 5/5
runs and \texttt{env} in 4/5. Beelzebub
\texttt{state\_cat\_marker}, \texttt{state\_echo\_marker}, and
\texttt{state\_ls\_tmp\_3} each timed out in 3/5 runs. DECEIVE
\texttt{cpuinfo} and \texttt{ps\_aux} timed out in runs 1 and 4. The complete
timeout accounting is given in Table~\ref{tab:timeouts}.

\begin{table}[ht]
\centering
\caption{Scored SSH timeout detail. The Beelzebub \texttt{exit} row is an unscored control row.}
\label{tab:timeouts}
\scriptsize
\begin{tabularx}{\textwidth}{l l r l}
\toprule
Target & Command & Runs & Run numbers \\
\midrule
Beelzebub & \texttt{exit} (control) & 1 & 1 \\
Beelzebub & \texttt{id} & 1 & 4 \\
Beelzebub & \texttt{os\_release} & 1 & 4 \\
Beelzebub & \texttt{passwd} & 1 & 5 \\
Beelzebub & \texttt{state\_cat\_marker} & 3 & 1,3,5 \\
Beelzebub & \texttt{state\_echo\_marker} & 3 & 1,4,5 \\
Beelzebub & \texttt{state\_ls\_tmp\_3} & 3 & 1,4,5 \\
Beelzebub & \texttt{state\_rm\_marker} & 2 & 2,3 \\
Beelzebub & \texttt{state\_rmdir\_bm\_dir} & 1 & 1 \\
Beelzebub & \texttt{whoami} & 1 & 2 \\
DECEIVE & \texttt{cpuinfo} & 2 & 1,4 \\
DECEIVE & \texttt{perm\_sudo\_l} & 1 & 3 \\
DECEIVE & \texttt{ps\_aux} & 2 & 1,4 \\
VelLMes & \texttt{env} & 4 & 1,2,4,5 \\
VelLMes & \texttt{ps\_aux} & 5 & 1,2,3,4,5 \\
\bottomrule
\end{tabularx}
\end{table}

The Beelzebub timeout rows total 17 because the table includes the run 1
\texttt{exit} control row. The scored total is 16 after excluding
\texttt{exit}; DECEIVE totals 5 and VelLMes totals 9. The detector's injected
strings were visible after ANSI stripping in all five runs for all 11
monitored command IDs.

\subsection{HTTP results}

Galah produced 11/100 5xx responses, six errors including three timeouts, six
distinct \texttt{Server} header values, 12 header-name-set inconsistencies,
0.2364 repeated-\texttt{GET /} similarity, 25/94 status agreements with nginx,
five documentation-example secret hits, and 19 target-port mentions. Its
latency median/p95 was 3.531/8.887 seconds. Ground-truth nginx produced zero
5xx responses, a constant \texttt{Server} header, 1.0000 repeat similarity,
and a 0.015-second median.

\subsection{HoneyBee results}

HoneyBee was evaluated with \texttt{glm-5p3-flash}, temperature 0.2, maximum
8192 tokens, and three generations per generator/case cell, for 18
generations. Docker Compose validity was 6/6 cases, Dockerfile validity 4/6,
and Nuclei validity 3/6. Length-limited results were 5/18, all in the
PostgreSQL case. Grafana was clean in 9/9 generations and PostgreSQL in 4/9.
The scorer was revised after the first results; original outputs remain in the
recorded validity file. These are generator-validity observations, not a
security verdict.

\section{Findings}

The following statements are limited to the supported findings in the
verified benchmark report:

\begin{itemize}[leftmargin=*]
\item Measured determinism differed between Cowrie (0.9874) and the LLM SSH
targets (0.4664--0.8364).
\item Galah produced six \texttt{Server} header values versus one for nginx,
0.2364 versus 1.0000 repeated-\texttt{GET /} similarity, and 25/94 status
agreement with nginx.
\item Beelzebub produced nine oversized outputs and 14/180 transcript-leakage
outputs; DECEIVE produced two oversized outputs, with a maximum of 135,108
characters.
\item Command-specific recurrence was visible in the timeout table: VelLMes
\texttt{ps\_aux} 5/5 and \texttt{env} 4/5; three Beelzebub state commands
3/5 each; and DECEIVE \texttt{cpuinfo} and \texttt{ps\_aux} in runs 1 and 4.
\item Median LLM SSH latency was reported as 3.0x for VelLMes, 3.4x for
AdvancedshelLM, 3.5x for DECEIVE, and 8.9x for Beelzebub relative to the
0.562-second ground-truth median. These ratios inherit the completed-command
latency rule.
\item The llm-honeypot injected strings were visible after ANSI stripping.
\end{itemize}

\section{Data-quality notes}

The harness uses a duplicate guard to skip complete target/run pairs and a
\texttt{--force} option that backs up and reruns an existing pair. The verified
raw data retain per-command timeout fields. Beelzebub's raw timeout total is
17 while its scored total is 16 because run 1 \texttt{exit} is a control row.

Large LLM fields motivated an increased CSV field-size limit. The analyzer
also bounds every similarity comparison after ANSI stripping to 4,000
characters, preventing very long outputs from dominating runtime. The final
data include outputs longer than 20,000 characters and a maximum Beelzebub
output of 239,959 characters.

HoneyBee scoring was revised after initial results; both the revised findings
and the existence of original outputs are recorded. The pilot and final
batches are not mixed. In the pilot summary, Galah has 29/100 5xx responses,
Beelzebub has zero timeouts, and DECEIVE has 13 timeouts. Available pilot raw
JSONL contains 14 DECEIVE timeouts, while
\texttt{results\_pilot/ssh\_commands\_long.csv} contains 17. The available
files therefore report different totals; no cause is assigned.

\section{Limitations}

The benchmark has five runs, one shared model/API endpoint, scripted inputs,
freshly started containers, censored timeout latencies, and a
character-similarity proxy. Personas, hostnames, and filesystem fixtures
differ between targets. Shared API load can affect timeout and latency
patterns. The 26-command fidelity intersection is a restricted subset, and
host-specific commands are reference-only. HoneyBee has three generations per
cell and static validity checks. These constraints limit generalization.

The overlapping fidelity confidence intervals do not support a tool ranking.
DECEIVE state-check failures do not establish its state-tracking capability.
The scripted observations do not establish attacker detectability, HoneyBee
output security, or behavior beyond this model, endpoint, and sample.

\section{Next steps and questions for discussion}

Several extensions follow directly from the documented limitations:

\begin{enumerate}[leftmargin=*]
\item Repeat the benchmark with a second model while preserving the same
command and request manifests.
\item Add an API-latency probe that separates service response time from
local transport and wrapper time.
\item Collect more runs and report confidence intervals with the same
completed-command and censoring rules.
\item Add a controlled state test with explicit fixtures and repeatable
checks.
\item Add detectability probes and a separate LLM-attacker evaluation.
\item Request real attacker session logs, if they can be obtained and handled
under appropriate controls.
\item Decide whether the next study should prioritize detectability, fidelity,
or LLM-attacker behavior; the present data do not select among these goals.
\end{enumerate}

\begin{thebibliography}{99}

% Bibliographic details must be completed before submission.
\bibitem{genai-sok}
SoK, \url{https://arxiv.org/abs/2510.25939}.

\bibitem{honeygpt}
HoneyGPT.

\bibitem{llmpot}
LLMPot, \url{https://arxiv.org/abs/2405.05999}.

\bibitem{galah}
galah, \url{https://github.com/0x4D31/galah}.

\bibitem{deceive}
DECEIVE, \url{https://github.com/coverflex-tech/DECEIVE}.

\bibitem{cowrie}
Cowrie, \url{https://github.com/cowrie/cowrie}.

\bibitem{beelzebub}
Beelzebub, \url{https://github.com/mariocandela/beelzebub}.

\bibitem{advancedshellm}
shelLM/AdvancedShelLM, \url{https://github.com/stratosphereips/AdvancedShelLM}.

\bibitem{vellmes}
VelLMes, \url{https://github.com/stratosphereips/VelLMes-AI-Deception-Framework}.

\end{thebibliography}

\end{document}
